\section{Riprogettazione}
\subsection{Osservazioni iniziali}
Per quanto riguarda la fase di riprogettazione abbiamo deciso di seguire un approccio rigoroso di seguito elencato:
\begin{enumerate}
	\item \textbf{Progettazione concettuale:} Creazione di mockup a bassa fedeltà per risolvere i problemi di usabilità riscontrati attraverso l'ausilio di pattern di progettazione della HCI
	\item \textbf{Progettazione fisica:} Layout, Elementi grafici, Colori e Testo
	\item Implementazione effettiva dei mockup nell'interfaccia 
\end{enumerate}

\subsection{Progettazione Concettuale}
\subsubsection{Software utilizzato}
Abbiamo deciso di usare un software chiamato "Excalidraw" che permette di fare mockup a bassa fedeltà direttamente nel browser. 
\subsubsection{Motivazioni}
Il motivo principale per cui abbiamo scelto di utilizzare mockup a bassa fedeltà rispetto a quelli in alta fedeltà risiede nel fatto che abbiamo preferito prima abbozzare il sistema per concentrarci esclusivamente sul risolvere i problemi di usabilità riscontrati senza perdere tempo/focus su altri aspetti meno critici
\clearpage
\subsubsection{Mockup}
In questa sezione saranno presentati i mockup che abbiamo realizzato con anche riferimento ai pattern che abbiamo utilizzato per risolvere tali problemi.
Ovviamente non siamo andati a sistemare tutti i problemi di usabilità riscontrati ma ci siamo concentrati maggiormente su quelli più critici.
Andremo a utilizzare il seguente schema:\\
\tabellaMockup{
	123
}{
	posizione example 
}{
	problema example 
}{
	descrizione example
}{
	\begin{itemize}
		\item \textbf{Pattern 1}
		\item \textbf{Pattern 2} 
	\end{itemize} 
}{
	\includegraphics[width=0.85\textwidth]{example_mockup_style} 
}

\subsubsection{Pattern HCI impiegati}
In questa sezione andremo a riepilogare i pattern che sono stati utilizzati con maggiori dettagli rispetto alla sezione precedente che si concentrava maggiormente sui mockup.

\subsection{Progettazione Fisica}
\subsubsection{Layout}
\subsubsection{Colori}
\subsubsection{Elementi grafici}
\subsubsection{Testo}

\subsection{Implementazione}
Di seguito verranno mostrate le sezioni riprogettate facendo riferimento ai relativi mockup, seguendo lo stesso schema a tabella visto per quest'ultimi.








